\chapter*{\chapterstyle{IV --- Intégrales Généralisées}}
\addcontentsline{toc}{section}{Intégrales Généralisées}
Aprés avoir défini précédemment l'intégrale de fonctions intégrables sur \textbf{sur un segment}, nous allons définir l'intégrale de fonction sur des intervalles \textbf{ouvert ou semi-ouvert}, ou encore les integrales de fonctions \textbf{non-bornées}, on les appellera alors \textbf{intégrales généralisées}.
\subsection*{\subsecstyle{Définition {:}}}
Soit \(I\) un intervalle \(\ioc{a}{b}\), et \(f\) une fonction intégrable sur \textbf{tout sous-segment} de \(\ioc{a}{b}\), alors on dira que \(f\) est \textbf{localement intégrable} et on définit alors une fonction auxiliaire \(\phi\) telle que :
\[
   \phi(x) = \int_{x}^{b} f(t) \d t
\]
Alors on peut définir l'intégrale généralisée, si elle existe, de \(f\) par:
\customBox{width=3cm}{
   \(\underset{x \rightarrow 0}{\lim} \; \phi(x)\)
}
Si cette limite existe, on dira alors que l'intégrale généralisée existe.\+ 
Dans le cas où \(f\) est \textbf{non-bornée} sur un intervalle \(I\), on peut alors séparer l'intégrale en deux sous-intégrales (ou plusieurs si les discontinuités sont multiples), qui pourront être étudiées de cette manière.\<

\underline{Exemple :} Etudions l'intégrale sur \(\ioc{0}{1}\) du logarithme. Cette fonction est bien localement intégrale sur l'intervalle, posons:
\[
   \phi(x) = \int_{x}^{1} \ln(t) \d t = - 1 - x\ln(x) + x 
\]
Or la limite quand \(x\) tends vers \(0\) existe et vaut \(-1\), et donc le logarithme est intégrable sur \(\ioc{0}{1}\) et on a \(\int_{x}^{1} \ln(t) \d t = -1\).
\subsection*{\subsecstyle{Intégrales de référence{:}}}
Dans cette section, on considère des intégrales sur un intervalle de \(\R_+\) ayant soit une borne supérieure infinie, soit un borne inférieure en 0.\+
Il existe alors plusieurs intégrales généralisées de référence, une des plus importante étant \textbf{l'intégrale de Riemann}:
\[
   \int \frac{1}{t^\alpha} \d t   
\]
Par une simple intégration directe suivie d'une disjonction de cas sur la limite de la primitive, on en déduit des propriétés importantes pour la suite:
\begin{align*}
   &\bullet \;\; \text{Elle \textbf{converge} en \(+\infty\) si et seulement si \(\alpha > 1\).} \\
   &\bullet \;\; \text{Elle \textbf{converge} en \(0\) si et seulement si \(\alpha < 1\).} \\
   &\bullet \;\; \text{Si \(\alpha = 1\), elle \textbf{diverge toujours}}
\end{align*}

Une autre intégrale de référence est \textbf{l'intégrale de Bertrand}:
\[
   \int \frac{1}{t^\alpha \ln(t)^\beta} \d t   
\]
Par comparaison avec l'intégrale de Riemann ci-dessus, on peut alors en déduire les propriétés suivantes:
\begin{align*}
   &\bullet \;\; \text{Elle \textbf{converge} en \(+\infty\) si et seulement si \(\alpha > 1\).} \\
   &\bullet \;\; \text{Elle \textbf{converge} en \(0\) si et seulement si \(\alpha < 1\).} \\
   &\bullet \;\; \text{Si \(\alpha = 1\), elle \textbf{converge} si et seulement si \(\beta > 1\).}
\end{align*}
\subsection*{\subsecstyle{Théorèmes de comparaison {:}}}
Soit \(f, g\) des fonctions \textbf{positives} sur l'intervalle considéré telles que \(f \leq g\), alors on a la propriété\footnote[1]{Ce cas découle directement de la positivité des fonctions considérées, en effet si \(f\) est positive, l'intégrale de \(f\) est \textbf{croissante}.} suivante:
\begin{align*}
   &\bullet \;\; \text{Si l'intégrale de \(g\) converge, alors l'integrale de \(f\) converge.} \\
   &\bullet \;\; \text{Si l'intégrale de \(f\) diverge, alors l'integrale de \(g\) diverge.}
\end{align*}
En particulier, cette propriété s'étend aux relations de négligeabilité\footnote[2]{Découle du fait que ce sont des majorations à partir d'un certain rang.}, en effet si \(f = o(g)\), alors:
\begin{align*}
   &\bullet \;\; \text{Si l'intégrale de \(g\) converge, alors l'integrale de \(f\) converge.} \\
   &\bullet \;\; \text{Si l'intégrale de \(f\) diverge, alors l'integrale de \(g\) diverge.}
\end{align*}
Enfin, on peut aussi considérer le cas de deux fonction \(f, g\) \textbf{équivalentes}\footnote[3]{En particulier, ce théorème nous permet d'utiliser tout l'arsenal des développements limités.} au point de discontinuité, et alors on a le théorème fondamental suivant:
\customBox{width=9cm}{
   Les intégrales de \(f\) et \(g\) sont \textbf{de même nature}.
}
\begin{center}
   \textit{Munis de ces théorèmes de comparaison et des intégrales de référence, on peut alors étudier efficacement la convergence des intégrales de fonctions de signe constant.}
\end{center}
\subsection*{\subsecstyle{Absolue convergence {:}}}
Dans le cas de fonction qui ne sont pas à signe constant sur l'intervalle étudié, on peut alors montrer le théorème suivant:
\customBox{width=9cm}{
   \(\text{Si } \int |f(t)| \d t \text{ converge, alors }  \int f(t) \d t \text{ converge}\)
}
On dira alors que \(f\) est \textbf{absolument convergente}\footnote[4]{La preuve utilise le fait que \(|f(t)| = f^+ + f^-\) avec \(f^+(t) = \max(f(t), 0)\) et \(f^+(t) = \max(-f(t), 0)\)}.
\begin{center}
   \textit{Cette propriété nous permet donc d'étendre nos théorèmes de comparaison aux fonctions qui changent de signe sur l'intervalle.}
\end{center}
\subsection*{\subsecstyle{Fonctions définies par une intégrale {:}}}
On peut aussi définir des fonctions donc le l'image sera définie par une intégrale, et donc dont le domaine de défintion sera l'ensemble des points tels que \textbf{l'intégrale convergence}. Une telle fonction est de la forme:
\customBox{width=6cm}{
   \(f(x): x \mapsto \int_{a}^{b} g_x(t) \d t\)
}
\underline{Exemple:} On peut définir la \textbf{fonction gamma}, définie sur \(\R \backslash \Z_-\) telle que:
\[
   \Gamma(x): x \mapsto \int_{0}^{+\infty} e^{-t}t^{x-1} \d t   
\]


\chapter*{\chapterstyle{IV --- Séries Numériques}}
\addcontentsline{toc}{section}{Séries Numériques}
Soit \((u_n)_{n\in\N}\) une suite, on appelle \textbf{série numérique de terme général \(u_n\)} la suite \(S_n\) des \textbf{sommes partielles} des termes de \((u_n)_{n\in\N}\), plus formellement, on a:
\[
   S_n = \sum_{k=0}^{n} u_k   
\]
Ce chapitre consistera en l'étude de la convergence ou de la divergence de telles séries. 
\subsection*{\subsecstyle{Propriétés{:}}}
Il est tout d'abord important de remarquer que la fonction qui à une suite associe sa série est \textbf{linéaire}, aussi on peut remarquer une propriété élémentaire:
\[
   u_n = S_{n} - S_{n-1}   
\]
Cette propriété nous permet alors de montrer la propriété fondamentale suivante:
\customBox{width=10cm}{
   Si la série \textbf{converge} alors son terme général tends vers \(0\).
}
Enfin dans certains cas, on peut exprimer le terme général \(u_n\) sous la forme \(a_{n+1} - a_n\) pour \((a_n)\) une autre suite. On appelera une telle série \textbf{série téléscopique} et on a la simplification suivante par linéarité:
\[
   \sum_{k=0}^{n} u_k = \sum_{k=0}^{n} (a_{k+1} - a_k) = a_{n+1} - a_0   
\]
On peut aussi remarquer que les fonctions exponentielle et logarithme étant continues, on peut définir de manière analogue la \textbf{suite des produits partiels} d'une suite, qui se ramène alors par passage au logarithme à l'étude d'une série.
\subsection*{\subsecstyle{Séries de référence{:}}}
Le cas le plus élémentaire de série est le cas des \textbf{séries géométriques}:
\[
   S_n = \sum_{k=0}^{n} q^k
\]
\begin{center}
   \textit{Cette série converge si et seulement si sa raison est inférieure à 1 en valeur absolue.}
\end{center}

Un second exemple remarquable est celui des \textbf{séries de Riemann} qui est analogue aux intégrales du même nom:
\[
   S_n = \sum_{k=0}^{n} \frac{1}{k^\alpha}
\]
\begin{center}
   \textit{Cette série converge si et seulement si \(\alpha\) est strictement supérieur à \(1\).}
\end{center}

Enfin un dernier exemple remarquable est celui des \textbf{séries de Bertrand} qui est aussi analogue aux intégrales du même nom:
\[
   S_n = \sum_{k=0}^{n} \frac{1}{k^\alpha \ln(k)^\beta}
\]
\begin{center}
   \textit{Cette série converge si et seulement si \(\alpha > 1\) ou si \(\alpha = 1\) et \(\beta > 1\).}
\end{center}
\subsection*{\subsecstyle{Théorèmes de comparaison {:}}}
Soit \((u_n),(v_n)\) des suites \textbf{positives}, telles que \((u_n) \leq (v_n)\) alors de manière analogue aux intégrales généralisées, alors on a la propriété\footnote[1]{Ce cas découle directement de la positivité des suites considérées, en effet si \((u_n)\) est positive, sa série associée est \textbf{croissante}.} suivante:
\begin{align*}
   &\bullet \;\; \text{Si la série associée à \(v_n\) converge, alors la série associée à \(u_n\) converge.} \\
   &\bullet \;\; \text{Si la série associée à \(u_n\) diverge, alors la série associée à \(v_n\) diverge.}
\end{align*}
En particulier, cette propriété s'étend aux relations de négligeabilité\footnote[2]{Découle du fait que ce sont des majorations à partir d'un certain rang.}, en effet si \(u_n = o(v_n)\), alors:
\begin{align*}
   &\bullet \;\; \text{Si la série associée à \(v_n\) converge, alors la série associée à \(u_n\) converge.} \\
   &\bullet \;\; \text{Si la série associée à \(u_n\) diverge, alors la série associée à \(v_n\) diverge.}
\end{align*}
Enfin, on peut aussi considérer le cas de deux suites \((u_n), (v_n)\) \textbf{équivalentes}\footnote[3]{En particulier, ce théorème nous permet d'utiliser tout l'arsenal des développements limités.}, et alors on a le théorème fondamental suivant:
\customBox{width=10cm}{
   Les séries associées à \(u_n\) et \(v_n\) sont \textbf{de même nature}.
}
\begin{center}
   \textit{Munis de ces théorèmes de comparaison et des séries de référence, on peut alors étudier efficacement la convergence des séries dont le terme général est de signe constant.}
\end{center}
\subsection*{\subsecstyle{Absolue convergence {:}}}
Dans le cas de séries dont le terme général n'est pas à signe constant, on peut alors montrer le théorème suivant:
\customBox{width=9cm}{
   \(\text{Si } \sum_{k=0}^{n} |u_k| \text{ converge, alors }  \sum_{k=0}^{n} u_k \text{ converge}\)
}
On dira alors que \(S_n\) est \textbf{absolument convergente}\footnote[4]{La preuve utilise le fait que \(|(u_k)| = (u_k)^+ + (u_k)^-\) avec \((u_k)^+ = \max(u_k, 0)\) et \((u_k)^+ = \max(-u_k, 0)\)}.
\begin{center}
   \textit{Cette propriété nous permet donc d'étendre nos théorèmes de comparaison aux séries dont le terme général change de signe.}
\end{center}
\subsection*{\subsecstyle{Critères spécifiques {:}}}
Soit \((u_n)\) une suite à termes \textbf{non-nuls}, alors si \(|\frac{u_{n+1}}{u_n}|\) admet une limite \(l\), alors\footnote[5]{\label{1}La preuve s'obtient en comparant \(u_n\) à une série géométrique} on a la \textbf{règle de d'Alembert}:
\begin{align*}
   &\bullet \;\; \text{Si \(l < 1\), la série associée converge.} \\
   &\bullet \;\; \text{Si \(l > 1\), la série associée diverge.} \\
   &\bullet \;\; \text{Sinon on ne peut pas conclure.}
\end{align*}
Soit \((u_n)\) une suite quelconque, alors si \(\sqrt[n]{|u_{n}|}\) admet une limite \(l\), alors\footref{1} on a la \textbf{règle de Cauchy}:
\begin{align*}
   &\bullet \;\; \text{Si \(l < 1\), la série associée converge.} \\
   &\bullet \;\; \text{Si \(l > 1\), la série associée diverge.} \\
   &\bullet \;\; \text{Sinon on ne peut pas conclure.}
\end{align*}
\begin{center}
   \textit{Ces règles spécifiques permettent d'obtenir la nature d'une série en étudiant simplement le terme général.}
\end{center}

\pagebreak
\subsection*{\subsecstyle{Familles sommables {:}}}
On s'intéresse maintenant à la possiblité d'étendre les propriétés de l'addition et de la multiplication aux sommes infinies. En particulier, est-il possible de donner une condition pour que la somme d'une série soit \textbf{invariante par permutation des termes}, \textbf{associative}, ou encore que le \textbf{produit de Cauchy} de deux séries convergentes soit convergente ?\<

On peut répondre affirmativement à ces question en introduisant le concept de \textbf{famille sommable}, et on peut alors montrer\footnote[6]{Hautement non trivial ...} que si la série associée à \(u_n\) est \textbf{absolument convergente}, alors:
\begin{align*}
   &\bullet \;\; \text{L'ordre des termes dans la somme ne change pas la somme} \\
   &\bullet \;\; \text{L'ordre des termes dans la somme ne change pas la somme}
\end{align*}

\chapter*{\chapterstyle{IV --- Suites de fonctions}}
\addcontentsline{toc}{section}{Suites de fonctions}
Soit \(E \subseteq F\), on appelle \textbf{suite de fonctions} une suite \((f_n)_n\in\N\) dont les éléments sont des fonctions de \(E\) dans \(F\). Dans ce chapitre, on cherchera à définir une notion de convergence de telle suites.

\subsection*{\subsecstyle{Convergence simple{:}}}
Une première approche serait de définir la convergence de la suite \((f_n)\) vers une fonction limite \(f\) comme:
\[
   \forall x \in \R \; , \; \forall \epsilon > 0 \; , \; \exists N \in \N \; , \; \forall n > N \; ; \; |f_n(x) - f(x)| < \epsilon
\]
Cela signifirait que pour un \(x_0\) fixé, la suite de \textbf{réels} \((f_n(x_0))\) converge vers \(f(x)\).\<

Par exemple si on considère la suite de fonctions définie par \(f_n(x) = x^n\) sur \(\icc{0}{1}\) et qu'on fixe \(x_0 \in \icc{0}{1}\) on a:

\begin{center}
   \begin{tikzpicture}[domain=0:1, scale=3.25, xscale=2.5]
      \draw[-latex] (0,0) -- (1,0) node [right] {$x$};
      \draw[-latex] (0, 0) -- (0,1) node [above] {$y$};
      \foreach \n in {1,...,15}
      {
         \draw[color=DarkBlue1] plot[smooth] (\x, {\x^\n});
         \fill[color=BrightRed1] (0.75, 0.75^\n) circle[radius=0.35pt, xscale=0.5];  
      }
      \draw[dashed, color = BrightRed1] (0.75, 0) -- (0.75,1);
      \draw[color=BrightRed1] node at (0.75, -0.1) {$x_0$}; 
   \end{tikzpicture}
\end{center}

Alors cette suite converge vers la fonction:
\[
   \begin{aligned}
      f : \icc{0}{1} &\longrightarrow \R\\
      x &\longmapsto \begin{cases}
         0 \text{ si } x < 1 \\
         1 \text{ sinon}
      \end{cases} 
   \end{aligned}
\]
On remarque un premier problème venant de cette définition, en effet les fonctions \(f_n\) sont toutes continues mais la limite n'est plus continue. Cela motive une notion de convergence plus \textbf{forte}.
\subsection*{\subsecstyle{Convergence uniforme{:}}}
On dit que \((f_n)\) \textbf{converge uniformément} vers \(f\) sur \(\R\) et on note \((f_n) \rightrightarrows f\) si et seulement si:
\[
   \forall \epsilon > 0 \; , \; \exists N \in \N \; , \; \forall n > N \; , \; \forall x \in \R \; ; \; |f_n(x) - f(x)| < \epsilon
\]
L'inversion des quantificateurs nous donne alors un unique seuil à partir duquel \(|f_n(x) - f(x)| < \epsilon\), en particulier, on peut alors montrer que cette notion de convergence équivaut à un criète plus pratique. En effet \((f_n)\) converge uniformément vers \(f\) si et seulement si:
\customBox{width = 4cm}{
   \(
      \vectNorm{f_n - f}_\infty^\R \rightarrow 0
   \)
}
On rapelle que la norme infini est définie par:
\[
   \vectNorm{f_n - f}_\infty^\R = \underset{x \in \R}{\sup}\left\{\left|f_n(x) - f(x)\right|\right\}
\]
Géométriquement, les boules pour cette norme consistent en des "tunnels" autour de la fonction centrale, et la convergence pour cette norme implique donc qu'à partir d'un certain rang, la suite de fonction se trouve toujours dans un tunnel arbitrairement petit.
\pagebreak

Par exemple la suite de fonctions \(f_n(x) = \frac{n + 2}{n}\sin(x)\) converge uniformément vers \(f: x \mapsto \sin(x)\) et graphiquement on a:

\begin{center}
   \begin{tikzpicture}[
      >=stealth,scale=0.75,line cap=round,
      bullet/.style={circle,inner sep=0.5pt,fill}
   ]
      \draw[->] (0,-3) -- (0, 3);
      \draw[] (-0.1,-3) -- (0.1, -3) node at (-0.45, -3) {$-3$};
      \draw[] (-0.1,0) -- (0.1, 0) node at (-0.3, 0) {$0$};
      \draw[] (-0.1,1) -- (0.1, 1) node at (-0.3, 1) {$1$};
      \draw[] (-0.1,2) -- (0.1, 2) node at (-0.3, 2) {$2$};
      \draw[] (-0.1,-2) -- (0.1, -2) node at (-0.45, -2) {$-2$};
      \draw[] (-0.1,-1) -- (0.1, -1) node at (-0.45, -1) {$-1$};

      \draw[color=black, domain=0:12, smooth, line width=0.45mm]   plot (\x,{sin(\x r)}) node[right] {$\sin(t)$}; 
      \draw[color=black!70,name path=A, domain=0:12, smooth, dashed, line width=0.3mm]   plot (\x,{sin(\x r) + 1}); 
      \draw[color=black!70,name path=B, domain=0:12, smooth, dashed, line width=0.3mm]   plot (\x,{sin(\x r) - 1}); 
      \foreach \n in {1, 2, 4, 6, 8, 10, 12, 14, 16, 18, 20, 22, 24}
      {
         \draw[color=BrightRed1!95, domain=0:12] plot[smooth, tension=0.70] (\x, {(\n+2)/\n*sin(\x r)});
      }
   \end{tikzpicture}   
\end{center}
\subsection*{\subsecstyle{Propriétés de régularité conservées{:}}}
Soit \((f_n)\) une suite de fonction qui converge uniformément vers \(f\), la valeur de notre définition est mise en lumière par les propriétés suivantes:
\begin{itemize}
   \item Si chaque\(f_n\) est \textbf{bornée}, alors \(f\) est \textbf{bornée}.
   \item Si chaque\(f_n\) est \textbf{continue}, alors \(f\) est \textbf{continue}.   
   \item Si chaque\(f_n\) est \textbf{uniformément continue}, alors \(f\) est \textbf{uniformément continue}.
\end{itemize}
\begin{center}
   \textit{Beaucoup de propriétés de régularité passent à la limite uniforme.}
\end{center}
Néanmoins le cas de la dérivabilité est plus subtil, en effet considéront la suite \((f_n)\) définie par \(f_n(x) = x^{1+\frac{1}{2n+1}}\), alors elle est bien définie et dérivable\footnote[1]{Remarquer que \(2n+1\) est toujours impair, donc on prends des racines n-ièmes impaires.} sur \(\R\), mais on a:
\begin{center}
   \begin{tikzpicture}[
      >=stealth,xscale=1.75, yscale=1,line cap=round,
      bullet/.style={circle,inner sep=0.5pt,fill}
   ]
      \draw[->] (0,0) -- (0, 5);
      \draw[->] (-4,0) -- (4,0);
      \draw[] (-0.1,1) -- (0.1, 1) node at (-0.3, 1) {$1$};
      \draw[] (-0.1,2) -- (0.1, 2) node at (-0.3, 2) {$2$};      
      \draw[] (-0.1,3) -- (0.1, 3) node at (-0.3, 3) {$3$};
      \draw[] (-0.1,4) -- (0.1, 4) node at (-0.3, 4) {$4$};

      \draw[color=black, domain=-3:3, smooth, line width=0.5mm]   plot (\x,{abs(\x)}) node[right] {$|t|$}; 
      \foreach \n in {1,..., 10}
      {
         \draw[color=BrightRed1!95, domain=-3:3] plot[smooth, tension=0.70] (\x, {\x*\x^(1/(2*\n+1))});
      }
   \end{tikzpicture}   
\end{center}
\[
   (f_n) \rightrightarrows f \text{ pour } f : x \mapsto |x|
\]
Et donc en particulier on a exhibé une suite de fonctions dérivables donc la limite uniforme \textbf{n'est pas dérivable} en un point.
\subsection*{\subsecstyle{Limites en un point{:}}}
On peut généraliser le résultat sur la continuité à un résultat sur les limites, en effet si \((f_n)\) est une suite de fonctions qui converge uniformément vers \(f\) et telles que \(f_n(x)\) admette une limite un un point \(a\), alors on a \textbf{le théorème d'interversion} suivant:
\customBox{width=7.5cm}{
   \[
      \lim_{n \rightarrow +\infty}\left(\lim_{x \rightarrow a} f_n(x)\right) = \lim_{x \rightarrow a}\left(\lim_{n \rightarrow +\infty} f_n(x) \right)
   \]
}
\subsection*{\subsecstyle{Intégrabilité{:}}}
On peut alors aussi montrer que si \((f_n)\) est une suite de fonction intégrables qui convergent uniformément vers \(f\), alors \(f\) est intégrable et on a \textbf{le théorème d'interversion} suivant:
\customBox{width=7cm}{
   \[
      \lim_{n \rightarrow +\infty}\int_{a}^{b} f_n(t) \d t = \int_{a}^{b} \lim_{n \rightarrow +\infty} f_n(t) \d t
   \]
}
\subsection*{\subsecstyle{Retour sur la dérivabilité{:}}}
Aprés avoir montré que la dérivabilité ne passe pas à la limite uniforme, on veut néanmoins trouver des critères permettant de trouver la dérivée d'une limite de fonctions, on considère alors une suite de fonctions \((f_n)\) toutes dérivables, et on suppose en outre que:
\begin{itemize}
   \item La suite de fonctions converge simplement pour au moins un \(x_0\).
   \item La suite des dérivées \((f_n')\) converge uniformément.
\end{itemize}
Alors \((f_n)\) converge uniformément vers une fonction dérivable \(f\) et:
\[
   f'(x) = \lim_{n \rightarrow +\infty} f_n'(x)
\] 
On note alors que pour que la dérivabilité passe à la limite, il nous faut la convergence uniforme de la suite des dérivées \textbf{et} la convergence de la suite en au moins un point.
\subsection*{\subsecstyle{Approximations uniformes{:}}}
On se demande maintenant quand peut-on dire qu'une fonction continue est limite uniforme d'une suite de fonctions. Le \textbf{théorème de Weierstrass} nous donne alors une réponse forte:
\begin{center}
   \textbf{Toute fonction continue sur un segment est limite uniforme d'une suite de fonctions polynomiales.}
\end{center}

\chapter*{\chapterstyle{IV --- Séries de fonctions}}
\addcontentsline{toc}{section}{Séries de fonctions}
On appelle série de fonctions de terme général \(f_n\) la suite \(S_n\) des \textbf{sommes partielles} des termes de la suite de fonctions \((f_n)_{n\in\N}\), plus formellement, on a:
\[
   S_n = \sum_{k=0}^{n} f_k   
\]
Ce chapitre consistera en l'étude de la convergence simple et/ou uniforme de telles séries. 
\subsection*{\subsecstyle{Convergence simple{:}}}
\subsection*{\subsecstyle{Convergence uniforme{:}}}
\subsection*{\subsecstyle{Propriétés de régularité conservées{:}}}
\subsection*{\subsecstyle{Convergence normale{:}}}

\chapter*{\chapterstyle{IV --- Séries Entières}}
\addcontentsline{toc}{section}{Séries Entières}
On appelle série entières de terme général \(a_nz^n\) une série de fonctions \(S_n\) de variable complexe \(z\) de la forme:
\[
   S_n = \sum_{k=0}^{n} a_nz^n 
\]
\subsection*{\subsecstyle{Rayon de convergence{:}}}
\subsection*{\subsecstyle{Convergence uniforme{:}}}
\subsection*{\subsecstyle{Propriétés de régularité conservées{:}}}
\subsection*{\subsecstyle{Convergence normale{:}}}

\chapter*{\chapterstyle{IV --- Séries de Fourier}}
\addcontentsline{toc}{section}{Séries de Fourier}